\section{Apollo and Dionysus}

\begin{quotex}
There are two kinds of people in the world: those who believe there are two kinds of people in the world and those who don't.
\end{quotex}


I've put off discussion of the final part of \textbf{Julius Evola}'s The Individual and the Becoming of the World, primarily because it is difficult for an \textbf{Apollonian} to critique the \textbf{Dionysian}. The latter wants to assert what the former prefers not to assert, and to deny what he would rather not deny. For example, \textbf{Empedocles} claims there are two forces: Love and Strife. Yet to admit this is tantamount to admitting Strife alone as the only force. The Apollonian wants to unite, harmonize, and reach the depths of agreement; hence, he can only admit the first force which has no opposite. Nevertheless, Evola brings a sword, not peace.

Now ``is'' can be used in two ways: it can refer to essence or it can refer to existence (as described in Essence and Existence\footnote{\url{http://www.gornahoor.net/?p=1917}}).

So when \textbf{Parmenides} claims that change is impossible, he is referring to the realm of unchanging ideas or essences. \textbf{Heraclitus} claims there is nothing but change and he means the realm of prime matter, appearances, or existence. Evola, following Aristotle and the Medievals, takes the middle path.

Now the Apollonian can only speak of ideas and concepts which follow the laws of logic. Existence, however, is inexplicable to him, it is alogical (not illogical). Essence follows Thought, Existence follows Will.

To make his point, Evola contrasts three fundamental attitudes, or spiritual tendencies, that separate the religious attitude (Apollo) from the initiatic attitude (Dionysus). Yet how is this the end of the story? Beyond yin and yang is the Tao. In this early work, there is no mention of the Tao, or the alchemical marriage of the opposites. Also, Evola adopts the Lunar side, the path of darkness, power (Shakti), irrationality, intoxication and ecstasy. For easy reference, these attitudes are shown in chart \ref{tab:ReligiousAttitudes}.

\begin{table}[h]\small\centering
\begin{tabular}{p{.5\textwidth}p{.5\textwidth}}\toprule
Religious Attitude & Attitude of the Mysteries \\\midrule

It believes that the world is fundamentally right from a principle of order, harmony, and goodness, and that everything opposed to it that is darkness, irrationality, indeterminateness, evil, either is illusion or something that must not be. & He faces up to reality, without veils, in its tragic, irrational, unprovidential nature \\\midrule

Consequently, it pushes such a principle back onto something ``other than me'' --- to something transcendent. We say ``consequently'' since my actual experience is showing me a world totally different from one that is ordered and rational, I am not entitled to also admit the actual existence, alongside it, of a providential and rational world, unless I relate it to something that transcends me and to believe that this something exists. & He rejects the belief in the present existence of an optimistic order, and the necessity of postulating the transcendental, the God to guarantee it, and to explain how it can coexist with the order of existing things that contradict it; \\\midrule 

If such a transcendent principle is intended to be the source of every reality, it must be said that the individual in himself is nothing and can be nothing. The ``I'' is annulled and its being is pushed back onto the other. & Thus the gulf remains between titanic pessimism and Dionysianism, bridged by recognizing or not recognizing the moral right to existence of a regulated, just, and rational world. In the case of the affirmative, he asserts the concept of the individual as a sufficient power in himself, capable of saving not only himself but also everything that he is connected to and that, at bottom, is the reflexion of what he concretely is. \\\bottomrule
\end{tabular}
\label{tab:ReligiousAttitudes}
\caption{Religious Attitudes}
\end{table}

\paragraph{Order and Harmony}


Apollo sees the world as ordered, harmonious, and providential. Whatever does not appear as such, is considered unreal or illusory. Dionysus, on the other hand, accepts the irrationality and tragedy of the world. Yes, but so do insects. For men, however, there have always been lawgivers, those who reveal the harmony hidden in the appearances and show the way to discover it. Even for Evola, the strong man is the one who can impose his will, as law. That the world is irrational and chaotic is the result of privation\footnote{\url{http://www.gornahoor.net/?p=1815}}, that is, something that should not be.

\paragraph{Transcendence}


Apollo notices that disharmony of existence. However, he does not thereby postulate a parallel existence of a rational and ordered world. Just the opposite. Unless he knows what is rational and harmonious, by what standard can he judge existence not to be so? Hence, Dionysus cannot even know whether existence is rational or irrational, without a transcendent standard. Hence, he must be his own lawgiver.

\paragraph{The Individual}


If there is a transcendent order, then the individual himself is part of that order, and cannot be said to exist. For Dionysus, ``man \emph{himself} has a real possibility of ordering and controlling this chaos and thereby bringing about a cosmos that doesn't already exist but demands its existence from him.'' This gets to the crux of Evola's point: the I is real, he orders his life, and the I knows himself from his life, which is his own reflexion.

\paragraph{The Royal Path}


This is the beginning, then, the initiation into the Mysteries. For Evola, and for Dionysus, the Self is not known as a concept or an object, but rather knows himself as the subject, the author of his life, and indirectly through viewing the world as his reflexion. The knowledge of the Mysteries resides in the darkness, beyond rational thought. That is why Dionysus is the god of wine: In vino, veritas, not in clear-eyed logic. Dionysus is the god of ecstasy, yet cannot give a logical response when asked about it. That is not the point in any case and to insist on it is to destroy the magic. Evola makes this contrast:

\begin{quotex}
To dare to tear away the veils with which Apollo hid primordial reality, to dare to transcend form in order to put oneself in contact with the primordial ``atrocity'' of a world in which good and evil, divine and human, rational and irrational, just and unjust no longer have any sense, being \emph{only} power, naked, free, fiery power; to dare that, and not to be swept away from this bottomless precipice, but to be able to do it, to control in oneself and, not surpassed but surpassing, to realize the wild pleasure of existing tragically --- such is the test of Dionysus, from which every will that truly wants to escape from the ``God of the Earth'' must have his consecration.
\end{quotex}



\flrightit{Posted on 2012-02-26 by Cologero}

\begin{center}* * *\end{center}

\begin{footnotesize}
\begin{sffamily}

\texttt{CaseyAnn on 2013-09-07 at 14:50 said:}

Cologero, I want to tell you how much I appreciate what you do here. \emph{Gornahoor} is the most essential and greatest website there could possibly be.

Since I'm in university, I don't have much time for Traditional readings on the side, so \emph{Gornahoor} is an especially excellent resource at this time. It's amazing that while I study philosophy in a profane academic setting I can simply search for the philosophers here and actually gain something truly valuable, see an understanding of Truth from what they have left us. Thank you.


\hfill

\end{sffamily}
\end{footnotesize}

